\documentclass[12pt]{article}
\usepackage[utf8]{inputenc}
\usepackage[T1]{fontenc}
\usepackage[margin=2.5cm]{geometry}
\usepackage{amsmath}

\begin{document}

\begin{flushright}
\today
\end{flushright}

\vspace{10pt}

\noindent Dear Editor,

\vspace{10pt}

\noindent We submit our manuscript entitled {\bf ``Dark energy as an information capacity deficit: Derivation from quantum causal graph theory and agreement with DESI DR2''} for consideration as an Article in {\it Nature Physics}.

\vspace{10pt}

\noindent{\bf Why Nature Physics?}

This paper reports a conceptual breakthrough in our understanding of dark energy: it is not a fluid, field, or cosmological constant---it is an {\it information capacity deficit}. We show that dark energy dynamics can be {\it derived} from two information-theoretic axioms, and that the resulting model's $(w_0,w_a)$ prediction falls within the DESI DR2 $1\sigma$ observed contour ($\chi^2=0.65$, $\Delta\chi^2=56.9\simeq7.5\sigma$ over $\Lambda$CDM).

This is not curve-fitting. Unlike all existing approaches---CPL, quintessence, emergent dark energy---which {\it parameterize} $w(z)$ and fit it to data, we derive $w(z)$ from the microphysics of quantum information modes being occupied by cosmic structure formation. The shape function $w(z)+1\propto{\rm SFR}(z)/[H(z)\rho_*(z)^{n/(n+1)}]$ is a {\it prediction}, with the single physical parameter $n=1$ (damping $\propto$ information deficit) selected by the data.

We believe this work meets {\it Nature Physics}' criteria of exceptional conceptual significance and broad interest:

\begin{enumerate}
\item {\bf Paradigm shift:} Dark energy is recast from a mysterious substance to a measurable consequence of the universe's finite information capacity. This connects cosmology to quantum information theory in a quantitatively testable way.

\item {\bf Timeliness:} The DESI DR2 results (2025--2026) have created intense interest in dynamical dark energy. Multiple {\it Nature Astronomy} papers have already appeared on the DESI signal. Our paper goes beyond phenomenological reporting to provide a {\it microphysical origin} for the observed dynamics.

\item {\bf Falsifiability:} DGF makes a ``smoking gun'' prediction that no other model can produce: $w(z)+1$ has a {\it peak} at $z\sim1$--$2$ due to the cosmic star formation history. DESI DR3 and Euclid can confirm or refute this within 2--3 years.

\item {\bf Cross-domain unification:} The same two axioms yield quantum mechanics (complementing Zilly 2026), the Einstein equations ($f(R)$ scalar-tensor), the black hole entropy area law (combinatorial proof via min-cut), and dark energy dynamics. The paper focuses on the dark energy result; the broader framework is documented in the Supplementary Information.

\item {\bf Laboratory testability:} We propose a specific quantum Darwinism experiment ($\sim$20 superconducting qubits) that can independently verify the $q$-field's existence, giving it an operational definition outside cosmology.
\end{enumerate}

\vspace{10pt}

\noindent{\bf Fit within existing literature}

The DESI Collaboration's DR2 results (arXiv:2503.14738, published in multiple journals) established the $2.8$--$4.2\sigma$ preference for dynamical dark energy. Follow-up phenomenological studies ({\it Nature Astronomy} {\bf 9}, 1879, 2025; Capozziello et al.\ 2026; Zhang et al.\ 2026) have explored CPL, $w_0w_a$CDM, and quintessence fits to the DESI data. All of these treat dark energy as a phenomenon to be parameterized.

Our paper is fundamentally different: we derive $w(z)$ from physical principles, and then compare with data. The agreement is not imposed---it emerges from the physics, with the natural benchmark $n=1$ selected by the data ($<0.2\sigma$ from the $\chi^2$ minimum).

Prior theoretical attempts to connect gravity to information (Verlinde 2011, Jacobson 1995) recovered Newtonian limits or Einstein equations but did {\it not} predict dark energy dynamics. The Quantum Memory Matrix model (Neukart et al.\ 2025) shares the information-theoretic spirit but predicts $|w_0+1|\sim10^{-2}$, which is excluded by DESI at $>10\sigma$.

\vspace{10pt}

\noindent{\bf Manuscript details}

\begin{itemize}
\item Main text: $\sim$3,200 words, 3 display items (Fig.~1--3, Table~1)
\item Supplementary Information: 8 sections covering full derivations, error analysis, extended competitor comparison, and laboratory test proposal
\item Data and code: available at [repository URL]; machine-readable contour data package included
\item All authors have approved the manuscript and its submission to {\it Nature Physics}
\end{itemize}

\vspace{10pt}

\noindent{\bf Suggested reviewers}

\begin{enumerate}
\item Expert in dark energy theory and DESI data analysis (e.g., a DESI collaboration theorist)
\item Expert in quantum information theory and emergence of spacetime
\item Expert in entropic/information-theoretic gravity approaches
\item Expert in observational cosmology (BAO, SNe, CMB joint constraints)
\end{enumerate}

\vspace{10pt}

\noindent We confirm that this work is not under consideration elsewhere and that there are no competing financial interests.

\vspace{20pt}

\noindent Sincerely,\\
\vspace{10pt}
\noindent The Authors

\end{document}
